\section{Signed Hurwitz numbers and their elementary multiplicities}\label{p8-sec:characters}

Let $a,b$ be the theta characteristics of the two normal lines,
$w=ba^{-1}\in H^1(C;\ZZ/2)$ their difference, and
$\epsilon=(-1)^{h^0(a)+h^0(b)}$. Put $H_m=1+m(g-1)$ and
$c_m=L_{m,H_m}$. For a cell $u$ of a partition $\lambda$, let
$h(u)$ denote its hook length and $c(u)$ its column minus its row.
Write $H_\lambda=\prod_{u\in\lambda}h(u)$. Let $\mathcal P_w$
consist of all partitions if $w=0$ and all self-conjugate partitions
if $w\ne0$, including the empty partition in either case. Set
\[
 Z_w(y)=\sum_{\lambda\in\mathcal P_w}
               H_\lambda^{\,2g-2}y^{|\lambda|}.
\]

\begin{theorem}\label{p8-thm:hooks}
The connected local counts at the first two possible source genera
satisfy $c_m=\epsilon^m[y^m]\log Z_w(y)$ and
\begin{equation}
 L_{m,H_m+1}
 =-\frac{\epsilon^m(g-1)}{12}[y^m]\,
  \frac{\displaystyle\sum_{\lambda\in\mathcal P_w}
         H_\lambda^{\,2g-2}\!
         \left(\sum_{u\in\lambda}h(u)^2\right)y^{|\lambda|}}
       {Z_w(y)}.
 \label{p8-eq:character-first-excess}
\end{equation}
\end{theorem}

For equal spin structures the right side is the first correction
in the balanced hook formula. The restriction to self-conjugate
partitions comes from a sign on one handle of the target.

\subsubsection{The spin sign of a collision}
At a nodal cover let $f:Y\to C$ be its \'etale normalization,
including all components. At a smooth cover with one triple
ramification point let $f:Y\to C$ be the cover itself. Write
$R$ for the ramification divisor of $f$, so $R=0$ in the first
case and $R=2r$ in the second. The two signs in the boundary
intersection are the parities of
\[
       f^*a\otimes\mathcal O_Y(R/2),\qquad
       f^*b\otimes\mathcal O_Y(R/2).
\]
For an odd branch these are the spin structures of spin Hurwitz
theory \cite{eop}. Their product has a simpler expression.

\begin{lemma}\label{p8-lem:odd-branch-sign}
Let $\alpha\in H^1(C;\ZZ/2)$ be the sign character of the sheet
permutations of $f$. Then the product of the two spin signs equals
\begin{equation}
       \epsilon^m(-1)^{\langle w\smile\alpha,[C]\rangle}.
 \label{p8-eq:odd-branch-sign}
\end{equation}
\end{lemma}
\begin{proof}
The bundle $E=f_*\mathcal O_Y(R/2)$ has rank $m$ and carries the
orthogonal form $\operatorname{Tr}(uv)$. Around $x=z^3$ its ramified
summand has basis $z^{-1},1,z$ over $\CC\{x\}$. The matrix of the
trace form in this basis has antidiagonal entries $3$ and all
other entries zero. The remaining sheets have the ordinary
permutation form. Thus the pairing extends nondegenerately over the
branch value. On its complement the determinant character of $E$
is the permutation sign. A 3-cycle is even, so this character
extends across the missing point and $w_1(E)=\alpha$.

For an orthogonal bundle of rank $m$ over a spin curve the
mod-two index formula is
\[
 h^0(a\otimes E)\equiv
 m h^0(a)+q_a(w_1(E))+\langle w_2(E),[C]\rangle\pmod2,
\]
where $q_a(v)=h^0(a\otimes v)-h^0(a)\pmod2$
\cite{ati,baraglia}. By the projection formula the left side is
the parity on $Y$. Adding this identity for $a$ and $b$ cancels
$w_2(E)$. Since $q_b(v)-q_a(v)=\langle w\smile v,[C]\rangle$,
the result follows. The same calculation applies to a disconnected
source by direct sum.
\end{proof}

\subsubsection{One signed handle and one marked collision}
We remove the common factor $\epsilon^m$ in degree $m$ until the
end of the calculation. For $w\ne0$ choose a simple nonseparating
circle representing its Poincar\'e dual as the first generator in a
surface presentation, disjoint from the branch value. The remaining
sign in~\eqref{p8-eq:odd-branch-sign} is $\operatorname{sgn}(A_1)$
for its sheet permutation $A_1$. Take $\chi=\operatorname{sgn}$ in
this case and $\chi=1$ for $w=0$. For $G=S_m$ put
\[
 K_\chi=\sum_{A,B\in G}\chi(A)\,ABA^{-1}B^{-1}\in\CC[G].
\]
On an irreducible representation $V$ of dimension $d$ and character
$\psi$, this central element has scalar
\[
 \frac{|G|}{d^2}\sum_{A\in G}\chi(A)|\psi(A)|^2
  =\left(\frac{|G|}{d}\right)^2
       {\bf1}_{\,V\otimes\chi\simeq V}.
\]
Indeed the sum over $B$ of the conjugates of $A^{-1}$ acts by
$|G|\overline{\psi(A)}/d$, and the last equality is character
orthogonality. The other $g-1$ handles contribute $K_1$.

The identity coefficient of $K_\chi K_1^{g-1}$ counts the labelled
\'etale covers with this sign. It equals the regular trace divided
by $|G|$. Relabelling sheets divides by $|G|$ once more, giving
the inverse automorphism weights. The summand of $V$ is therefore
$(|G|/d)^{2g-2}$ when $V\otimes\chi\simeq V$. For
$V=V_\lambda$ the hook-length formula gives $|G|/d=H_\lambda$,
and tensoring with sign transposes $\lambda$. The resulting
disconnected series is $Z_w$. Decomposition into connected
components gives the formula for $c_m$, as in \cite[Appendix~F]{WueUnramified}.

Denote by $N_m$ the signed mass over one fixed target point of
\'etale covers with two distinct sheets identified to give a
connected nodal cover. Denote by $T_m$ the signed mass of connected
covers with a single 3-cycle over that point. Both masses still omit
$\epsilon^m$. Before imposing connectivity the first mark has
$\binom m2$ choices on each labelled fiber. After gluing, the marked
component is unique, and all its companions are unmarked \'etale
covers. Removing their generating function gives
\begin{equation}
    \sum_{m\ge1}N_m y^m=\frac{y^2 Z_w''(y)}{2Z_w(y)}.
 \label{p8-eq:character-nodal}
\end{equation}
In particular this includes a node joining two distinct components
of the normalization. For the triple mark the handle product is
multiplied by the class sum of all 3-cycles. The Jucys--Murphy
elements $J_j=\sum_{i<j}(ij)$ satisfy
\[
     \sum_j J_j^2=\binom m2\,1+\sum_{\tau\text{ a 3-cycle}}\tau .
\]
Indeed the nondiagonal terms of $J_j^2$ list the 3-cycles with
largest letter $j$, each once. In the Young basis the $J_j$ act
by the contents of the cells containing $j$ \cite{jucys}. Removing
the unmarked components as above now gives
\begin{equation}
 \sum_{m\ge1}T_m y^m
  =\frac{1}{Z_w(y)}\sum_{\lambda\in\mathcal P_w}
      H_\lambda^{\,2g-2}
      \left(\sum_{u\in\lambda}c(u)^2-\binom{|\lambda|}{2}\right)
      y^{|\lambda|}.
 \label{p8-eq:character-triple}
\end{equation}

\begin{proof}[Proof of Theorem~\ref{p8-thm:hooks}]
The two collision intersections give the finite contribution
$-\epsilon^m(g-1)(3N_m+T_m)/12$. The contracted elliptic tails add
$-m(g-1)c_m/12$. Since $y(\log Z_w)'=y Z_w'/Z_w$, the three
terms over the common denominator $Z_w$ have total insertion on
a partition of size $m$ equal to
\[
 3\binom m2+\sum_{u\in\lambda}c(u)^2-\binom m2+m
 =m^2+\sum_{u\in\lambda}c(u)^2
 =\sum_{u\in\lambda}h(u)^2 .
\]
The last equality is the hook-content square identity \cite{menon}.
Restoring $\epsilon^m$ proves~\eqref{p8-eq:character-first-excess}.
\end{proof}

\subsubsection{Degree three}
Write $s=2g-2$. In degree three the partitions $(3),(2,1),(1^3)$
have hook products $6,3,6$, respectively. Only $(2,1)$ is
self-conjugate. Equations~\eqref{p8-eq:character-nodal}
and~\eqref{p8-eq:character-triple}, including the lower coefficients in
the division by $Z_w$, give
\[
\begin{array}{c|cc}
 &N_3&T_3\\ \hline
 w=0   &6\cdot6^s+3\cdot3^s-2\cdot2^s &(4^g-1)3^s\\
 w\ne0 &3^{2g-1}&-3^s .
\end{array}
\]
Thus the finite contribution for equal spins is
$-(g-1)(22\cdot6^s+8\cdot3^s-6\cdot2^s)/12$, and for distinct
spins it is $-2\epsilon(g-1)3^s/3$. At this prime degree the
proper divisor one contributes $\epsilon$ times the finite degree
of the balanced genus-$g$ core. The degree-one finite contribution
at positive excess is zero. Cancellation of the elliptic tails
therefore gives, for $a\ne b$,
\begin{equation}
        n_{3,H_3+1}
       =\frac{\epsilon(g-1)}{6}
          \left(11\cdot6^{\,2g-2}-3\cdot2^{\,2g-2}\right).
 \label{p8-eq:cubic}
\end{equation}
For $g=2$ the ordinary collision masses are $235$ and $135$,
with finite contribution $-70$. For distinct spins they are $27$
and $-9$ before restoring $\epsilon$, with finite contribution
$-6\epsilon$. Their difference gives $n_{3,5}=64\epsilon$.
\subsection{Elementary exponents and unramified quotients}\label{p8-sec:exponents}

Let $F_h(q)$ be the balanced antidiagonal local factor of genus $h$ in
the rational universal algebra of complex one-cycles of virtual dimension
zero \cite{pardon,WueProduct}. We use the convention that $F_h$ has
constant term one. Its image in the algebra of almost-complex cycle
families is the elementary cluster of genus $h$ of Ionel and
Parker \cite{ip,WueProduct}. For the isolated spin germ, the characteristic
insertions at the compact cycles $m[C]$ have a unique balanced lift
\cite[Proposition~3.1]{WueUnramified}
\begin{equation}\label{p8-eq:lifted-product}
  \widetilde P=\prod_{d\ge1,\ h\ge H_d}F_h(q^d)^{n_{d,h}}.
\end{equation}
The exponents are integral and have finite support in $h$ at fixed
$d$. This uses geometric decomposition after compact realization and
the multiplicity-one detection theorem of \cite{WueProduct}. Along the
fixed-target cluster decompositions, connected Gromov--Witten degrees
add on disjoint targets. Thus the same exponents give the expansion
of the connected series in the elementary logarithms. The use of
fixed-target degrees in this expansion relies on that geometric
comparison. In particular no evaluation of arbitrary analytic
parameter families is required here.

Put $n_d=n_{d,H_d}$. In this notation the tail cancellation of
\eqref{p8-eq:tails-exponent} subtracts the equal-spin finite-map integrals
on unramified covers. We next express the subtraction in terms of
the individual maps to $C$.

\subsubsection{A count using the maximal unramified quotient}
Fix a target point $x$. Let $\mathcal B_m(x)$ denote the isomorphism
classes of the finite maps over the nodal and triple divisors with
collision value $x$. Such an object has one distinguished node or one
triple ramification point. Write $k(f)=1$ in the nodal case and
$k(f)=3$ in the triple case, and write $\tau(f)$ for the product of
the two spin parities in~\eqref{p8-eq:intro-divisor}. All automorphisms
in this section fix the target. Put $\Gamma=\pi_1(C)$ and regard
$H(f)=f_*\pi_1(Y)$ as a conjugacy class of subgroups. This is the factorization by the image of the fundamental group used
by Bryan and Pandharipande \cite[\S2.2, Proposition~2.3]{bp01}.
It uses
the connected nodal space $Y$ itself when $f$ has a node.
The connected cover of $C$ corresponding to $H(f)$ is finite:
$f$ lifts to it, and the lift is holomorphic locally on every
component. A nonconstant holomorphic map from a compact component
to a connected Riemann surface forces that surface to be compact.
Composition of degrees shows that $[\Gamma:H(f)]$ divides $m$.
This cover is the maximal unramified quotient of $f$, in the sense
that every factorization through an unramified connected cover
factors through it after making compatible choices of base points.
For finite-index $H\le\Gamma$ write $\sigma(H)$ for the determinant
sign of the spin pair on its unramified cover.

\begin{theorem}\label{p8-thm:transfer}
For every $m\ge1$,
\[
 n_{m,H_m+1}
  =\frac{g-1}{4}\!
     \sum_{f\in\mathcal B_m(x)}
     \frac{\sigma(H(f))-\tau(f)}
          {k(f)\,|\operatorname{Aut}_C(f)|}.
\]
In particular the difference of the exponents for two spin pairs
is obtained by taking the corresponding difference in these
boundary and unramified signs on the same covers.
\end{theorem}
\begin{proof}
By the orbit formula for the determinant signs \cite[Theorem~16.1]{WueUnramified}, there are
integers $b_K$ indexed by conjugacy classes of finite-index subgroups
such that
\[
 \sigma(H)=\sum_{[K]} b_K |(\Gamma/K)^H|,
 \qquad n_d=\sum_{[\Gamma:K]=d} b_K .
\]
In the second sum each conjugacy class is used once. The first formula
is the fixed-point count, or mark, of a unit in the completed Burnside
ring; only subgroups of index dividing $[\Gamma:H]$ contribute.
For a fixed cover $z:Z\to C$ corresponding to $K$ of index $d$,
the number of lifts of $f$ to $Z$ is $|(\Gamma/K)^{H(f)}|$.
The groupoid of such lifts over an object $f$ has mass this number
divided by $|\operatorname{Aut}_C(f)|$, by orbit-stabilizer.
Here $Z$ is a fixed target, so its deck group does not divide this
number. Counting the same lifted maps on $Z$ allows the collision
value to be any of the $d$ preimages of $x$. The mass at each such
value is the same, by continuation along the connected curve $Z$.
An unramified change of target preserves the node or the triple
point and hence the weight $1/k(f)$. All boundary signs in the
equal-spin elementary pair on $Z$ are positive. By
Theorem~\ref{p8-thm:boundary} and $H_d-1=d(g-1)$ we obtain
\[
 F_{m/d}^{\rm elem}(H_d)=
 -\frac{g-1}{4}
  \sum_{f\in\mathcal B_m(x)}
  \frac{|(\Gamma/K)^{H(f)}|}
       {k(f)\,|\operatorname{Aut}_C(f)|}
\]
whenever $d\mid m$. For $d\nmid m$ the fixed sets are empty.
Multiply by $b_K$, sum, and apply~\eqref{p8-eq:tails-exponent} and
the signed boundary formula for $D$.
\end{proof}

\begin{corollary}\label{p8-cor:prime-growth}
Let $a\ne b$ and let $p$ be an odd prime. The signs of $n_{p,H_p}$
and $n_{p,H_p+1}$ are $-\epsilon$ and $\epsilon$, respectively, and
\[
 |n_{p,H_p+1}|\ > \frac{g-1}{4}\binom{p}{2}\,|n_{p,H_p}| .
\]
For each $g$ there is a path of tame structures on a closed
symplectic six-torus for which the isolated positive-sign core at
its two ends has different local elementary exponents at both these
source genera for every odd prime $p$.
\end{corollary}
\begin{proof}
Every collision in prime degree has $H(f)=\Gamma$. Otherwise its
maximal unramified quotient would have degree $p$ and the lift degree
one. A finite flat degree-one cover of a smooth curve is an
isomorphism, inconsistent with either collision. Thus its unramified
sign in Theorem~\ref{p8-thm:transfer} is $\epsilon$. The boundary sign
is $\epsilon^p$ times the permutation sign around a loop Poincar\'e
dual to the difference $ba^{-1}$, by the parity formula in
Section~\ref{p8-sec:characters}. Consequently each negative permutation
sign contributes with sign $\epsilon$. Let $U_-$ be the inverse-deck
mass of connected unramified covers with that negative sign. Each
one supplies $\binom p2$ unordered pairs in the fiber at $x$ whose
identification forms a node of the same character. Their total mass
is $\binom p2 U_-$. In least genus $n_{p,H_p}=-2\epsilon U_-$ by
the unramified formula. Also $U_->0$: choose surface generators
with the chosen loop first, send it to a transposition, send the
generator of another handle to a $p$-cycle, and send all remaining
generators to the identity. The source is connected and the surface
relation holds. These nodes alone give the weak inequality. There is also a nodal cover
with disconnected normalization of degrees $1$ and $p-1$. Choose a
connected cyclic $(p-1)$-cover whose sign on the chosen loop is negative
by sending the loop to a full cycle. Join one of its sheets to the
degree-one sheet at $x$. The resulting map has the same negative
permutation sign, giving a further positive mass after multiplication
by $\epsilon$. This proves the strict inequality.
The path and its primitive core class are supplied by the compact
spin realization of \cite[\S15]{WueUnramified}. Its ends are the equal
spin germ and a distinct-spin germ of the same positive sign. The
equal-spin germ has only the genus-$g$ elementary factor in degree
one. The displayed comparison concerns the endpoints; it requires
no assignment of a branched count to a single parameter of the path.
\end{proof}

One may express the result as a comparison of integral universal
products. Let $z$ be an indeterminate and let $Q_h(z)\in\ZZ[z]$ be the
polynomial with $Q_h(4\sin^2(t/2))=4\sin^2(ht/2)$. It is integral
by the cosine recurrence, with $Q_h(0)=0$ and $Q_1(z)=z$.
With $\mathcal P_w$ the partition set of \S\ref{p8-sec:characters} for the spin difference $w$, put
\[
 \mathcal Z(q,z)=\sum_{\lambda\in\mathcal P_w}
     (\epsilon q)^{|\lambda|}\prod_{s\in\lambda}Q_{h(s)}(z)^{g-1}.
\]
This has a unique formal expansion as an integral product of ordinary
balanced hook series of genera $h$ evaluated at $q^d$. Indeed, in
degree $d$ subtract the already determined factors of smaller degree
by multiplying by their inverses. The resulting coefficient is an
integral polynomial in $z$ divisible by $z^{d(g-1)}$. The degree-one
terms of the balanced factors are the distinct monomials $z^{h-1}$,
so its coefficients give the required integer exponents for this
degree. The divisibility is preserved under products and inversion,
and each degree requires only finitely many genera. Using these
exponents as in~\eqref{p8-eq:lifted-product} gives a balanced element
in the completed universal algebra.
\begin{corollary}\label{p8-cor:quotient}
This element agrees with $\widetilde P$ modulo the subgroup of
balanced products generated, with degree completion, by the
$F_h(q^d)$ with $h\ge H_d+2$.
\end{corollary}
\begin{proof}
Substitute $q=y t^{-2(g-1)}$ and $z=4\sin^2(t/2)$ in the formal
product. There are no negative powers of $t$ after this normalization.
The constant and quadratic coefficients in $t$ in its connected
logarithm are the unramified connected counts and
\eqref{p8-eq:character-first-excess}: a box of hook length $k$ contributes a
relative quadratic term $-(g-1)k^2/12$. These are exactly the
corresponding numerical coefficients of the geometric product.
Triangularity in covering degree and then in source genus, as in
\eqref{p8-eq:tails-exponent}, identifies its integer exponents in both
genera. Applying those exponents in the universal factors proves the
claim. Thus the two lowest genera can also be recovered directly from the collision and tail geometry, independently of the all-genus comparison in Theorem~\ref{thm:spin}.
\end{proof}
